\chapter{A compact guide to notation and mathematical judgment}\label{app:guide}
This appendix gathers in one place two things that the book develops gradually: its notation, and the mathematical judgment that guides how its proofs are read and found. Section~\ref{appE:sec:notation} is a table of the main symbols, each with the place where it is explained. Sections~\ref{appE:sec:structures}--\ref{appE:sec:record} collect the ways of reading and reasoning that run through the chapters: what the main structures of the book allow and what its main proofs actually use, how to choose the next move in a proof, how to turn a reading difficulty into a precise question, and how to keep a record of a proof you are studying. Section~\ref{appE:sec:rubric} suggests a rubric for the capstone investigation, and Section~\ref{appE:sec:finishing} describes what the course aims to leave you able to do.
\section{A notation map, with places to return to}\label{appE:sec:notation}
The table below lists the main symbols of the book, roughly in order of appearance, with the place where each is explained. A one-line summary is only a reminder: to use a symbol with confidence again, reread that section and its examples.

A few symbols, and many letters, do more than one job; the surrounding sentence and the kind of object involved decide which.
\begin{itemize}
\item Vertical bars give the absolute value $|x|$ of a real number or the size $|S|$ of a finite set; double bars, as in $\norm{x}_2$, give the length of a vector.
\item $(1,3)$ can be an interval or an ordered pair; $[1,3]$ is an interval, while $[x]$ is an equivalence class.
\item The arrow $\to$ appears in a function $f:A\to B$, a directed edge $u\to v$, and a convergent sequence $x_n\to x_*$; the barred arrow, as in $x\mapsto x^2$, names the rule of a function, and $\Rightarrow$ joins two statements.
\item When $f$ has an inverse, $f^{-1}$ is that inverse function, so $f^{-1}(y)$ is a single input, not the reciprocal $1/f(y)$ (Section~\ref{ch03:sec:functions}). The preimage $f^{-1}(T)$ of a set $T$ of outputs is a set of inputs, defined for every function; when $f$ has an inverse, it is the image of $T$ under the inverse (Section~\ref{ch03:sec:images}).
\item The letter $q$ is a quotient in Chapter~\ref{ch:induction} and a contraction factor in Chapters~\ref{ch:contraction} and~\ref{ch:capstone}. The capital $A$ is a set in Chapter~\ref{ch:maps}, a matrix in Chapter~\ref{ch:linear} and an event in Chapter~\ref{ch:probability}, while $A_s$ is an autocorrelation in Chapter~\ref{ch:signs}. The letter $I$ is an identity matrix in Chapter~\ref{ch:linear}, but $I_A$ is an indicator in Chapter~\ref{ch:probability}; and $e_j$ is a standard basis vector in Chapter~\ref{ch:linear}, but $e_n$ is the error at step~$n$ in Chapter~\ref{ch:capstone}.
\end{itemize}

\begin{longtable}{@{}p{1.25in}p{3.95in}@{}}
\toprule Symbol & Meaning, and where it is explained\\\midrule\endhead
$\NN,\ZZ,\QQ,\RR$ & The natural numbers $0,1,2,\ldots$ (zero included), the integers, the rational numbers, and the real numbers (Section~\ref{ch01:sec:notation}).\\
$\sum_{j=a}^{b}x_j$ & The sum $x_a+x_{a+1}+\cdots+x_b$, with one term for each integer $j$ from $a$ to $b$ (Section~\ref{ch01:sec:notation}). When $b<a$ there are no terms, and this \emph{empty sum} is $0$ (Section~\ref{ch05:sec:first-induction}). A sum such as $\sum_{v\in V}\deg(v)$ has one term for each member $v$ of the finite set $V$, added in any order (Section~\ref{ch07:sec:double}).\\
$x<y$, $x\le y$ & $x$ is strictly smaller than $y$; $x$ is smaller than or equal to $y$ (Section~\ref{ch01:sec:notation}).\\
$\sqrt t$ & For a real number $t\ge0$, the nonnegative number whose square is $t$; so $\sqrt9=3$, not $-3$ (Section~\ref{ch01:sec:notation}).\\
$P\Rightarrow Q$, $P\Leftrightarrow Q$ & ``If $P$, then $Q$''; ``$P$ if and only if $Q$,'' meaning both $P\Rightarrow Q$ and $Q\Rightarrow P$. Writing either states a claim, not that the claim has been proved (Section~\ref{ch02:sec:implication}).\\
$\neg,\land,\lor$ & Not; and; or. The ``or'' is inclusive: $P\lor Q$ is true when at least one of $P$ and $Q$ is true (Section~\ref{ch02:sec:connectives}).\\
$\forall,\exists$ & For every; there exists (at least one). Quantifiers are read from left to right: a value chosen at an $\exists$ may depend on the values chosen before it, but not on those chosen after it, so changing the order can change the meaning (Section~\ref{ch02:sec:quantifiers}).\\
$\eps,\delta$ & Traditional names for tolerances, positive amounts of error we are prepared to allow. They are names, not fixed numbers: ``for every $\eps>0$'' covers every positive value (Sections~\ref{ch02:sec:quantifiers} and~\ref{ch11:sec:distance}). In Chapter~\ref{ch:capstone}, $\eps_n$ is an update error (last row).\\
$|x|$ & The absolute value of a real number $x$: its distance from $0$ on the number line (Section~\ref{sec:magnitude}).\\
$[a,b]$, $(a,b)$, $[0,\infty)$ & Intervals of real numbers. A square bracket includes its endpoint and a round bracket excludes it; $[0,\infty)$ is the set of all real $x$ with $x\ge0$ (Section~\ref{sec:magnitude}).\\
$x\in A$, $x\notin A$ & $x$ is a member of the set $A$; $x$ is not a member of $A$ (Section~\ref{ch03:sec:sets}).\\
$A\subseteq B$, $\varnothing$ & Every member of $A$ is also a member of $B$; the empty set, which has no members (Section~\ref{ch03:sec:sets}).\\
$\{x\in A:P(x)\}$ & The set of members of $A$ that have the property $P$. Read the colon as ``such that'' (Section~\ref{ch03:sec:sets}).\\
$A\cup B$, $A\cap B$ & Union: the objects in at least one of $A$ and $B$. Intersection: the objects in both (Section~\ref{ch03:sec:sets}).\\
$A\setminus B$ & The set of members of $A$ that are not in $B$: members are removed, not numbers subtracted (Section~\ref{ch03:sec:sets}).\\
$A\times B$ & The set of ordered pairs $(a,b)$ with $a\in A$ and $b\in B$ (Section~\ref{ch03:sec:sets}).\\
$|S|$ & The number of members of a finite set $S$, each counted once; $|\varnothing|=0$ (Section~\ref{ch03:sec:sets}).\\
$f:A\to B$, $x\mapsto x^2$ & A function with domain $A$ and codomain $B$; the function that sends each input $x$ to $x^2$ (Section~\ref{ch03:sec:functions}).\\
$g\circ f$, $\id_A$, $f^{-1}$ & ``$g$ after $f$'': apply $f$ first, then $g$. The identity function on $A$, which sends every element to itself. The inverse of $f:A\to B$, when it exists: the function $f^{-1}:B\to A$ with $f^{-1}(f(x))=x$ for every $x\in A$ and $f(f^{-1}(y))=y$ for every $y\in B$ (Section~\ref{ch03:sec:functions}).\\
$f(S)$, $f^{-1}(T)$ & The image of a set $S$ of inputs: all outputs $f(x)$ with $x\in S$. The preimage of a set $T$ of outputs: all inputs $x$ with $f(x)\in T$ (Section~\ref{ch03:sec:images}).\\
$d\mid a$, $\gcd(a,b)$ & The positive integer $d$ divides the integer $a$: $a=dk$ for some integer $k$. The greatest common divisor of nonnegative integers $a$ and $b$, not both zero: the largest positive integer that divides both (Section~\ref{ch05:sec:euclid}).\\
$a\bmod m$ & The remainder when the integer $a$ is divided by the positive integer $m$; one of $0,1,\ldots,m-1$ (Section~\ref{ch05:sec:euclid}).\\
$x\sim y$, $[x]$ & $x$ is related to $y$; the equivalence class of $x$, the set of all elements related to $x$ (Section~\ref{ch06:sec:differences}).\\
$a\equiv b\pmod m$ & $a$ is congruent to $b$ modulo $m$: the positive integer $m$ divides $a-b$ (Section~\ref{ch06:sec:residues}).\\
$[a]_m$, $\ZZ/m\ZZ$ & The residue class of $a$: the whole set of integers congruent to $a$ modulo $m$, not just the one integer $a$. The set of all $m$ residue classes (Section~\ref{ch06:sec:residues}).\\
$x\preceq y$ & $x$ is below $y$ in a partial order, a relation that is reflexive, antisymmetric and transitive (Section~\ref{ch06:sec:order}).\\
$\deg(v)$ & The degree of a vertex $v$: the number of its neighbors (Section~\ref{ch07:sec:graph}).\\
$u\to v$ & A directed edge from $u$ to $v$; $v$ is then an out-neighbor of $u$ (Section~\ref{ch07:sec:sink}).\\
$N_1^+(v)$, $N_2^+(v)$ & The first and second out-neighborhoods of $v$: the vertices at directed distance exactly one, and exactly two, from $v$ (Section~\ref{ch07:sec:seymour}).\\
$N(s)$, $N(S)$ & The tasks accepted by the student $s$; the tasks accepted by at least one student of the set $S$ (Sections~\ref{ch08:sec:model} and~\ref{ch08:sec:shortage}).\\
$S\subsetneq L$ & $S$ is a proper subset of $L$: $S\subseteq L$ and $S\ne L$ (Section~\ref{ch09:sec:inductive}).\\
$\RR^n$ & The set of vectors $(x_1,\ldots,x_n)$ with $n$ real coordinates (Section~\ref{ch10:sec:vectors}).\\
$A_{ij}$, $Ax$ & The entry of the matrix $A$ in row $i$ and column $j$; the vector obtained by applying $A$ to $x$ (Section~\ref{ch10:sec:linearmaps}).\\
$\ker T$ & The kernel of a linear map $T$: the inputs that $T$ sends to the zero vector (Section~\ref{ch10:sec:kernel}).\\
$u\cdot v$ & The dot product: multiply corresponding coordinates and add. The result is a number, not a vector (Section~\ref{ch10:sec:inner}).\\
$\norm{x}_2$ & The Euclidean length $\sqrt{x\cdot x}$ of a vector. The subscript $2$ is part of the name, not a coordinate (Section~\ref{ch10:sec:inner}).\\
$A^{\mathsf T}$, $I$ & The transpose of $A$, obtained by exchanging rows and columns. The identity matrix, with ones on the diagonal and zeros elsewhere (Section~\ref{ch10:sec:product}).\\
$(X,d)$, $d(x,y)$ & A metric space: a set $X$ with a chosen metric $d$, which gives every pair of points $x,y$ a distance $d(x,y)$ and satisfies, for all points $x,y,z$: $d(x,y)\ge0$, with $d(x,y)=0$ exactly when $x=y$; $d(x,y)=d(y,x)$; and the triangle inequality $d(x,z)\le d(x,y)+d(y,z)$. The choice of metric is part of the problem (Section~\ref{ch11:sec:metric}).\\
$d_2$, $d_\infty$ & Two metrics on $\RR^n$: the Euclidean distance $\norm{x-y}_2$, and the largest of the coordinate differences $|x_j-y_j|$ (Section~\ref{ch11:sec:metric}).\\
$x_n\to x_*$ & The sequence converges to $x_*$: for every tolerance $\eps>0$, all terms from some index on are at distance less than $\eps$ from $x_*$ (Section~\ref{ch11:sec:convergence}).\\
$x_*$ & A name for a special point: a limit in Chapter~\ref{ch:limits}, a fixed point in Chapters~\ref{ch:contraction} and~\ref{ch:capstone}. The star is part of the name (Section~\ref{ch11:sec:convergence}).\\
$Tx$ & Short for $T(x)$ (Section~\ref{ch12:sec:iteration}).\\
$\Omega$, $p(\omega)$ & The sample space, a finite set of outcomes; the probability of the outcome $\omega$. The numbers $p(\omega)$ are nonnegative and add up to $1$ (Section~\ref{ch13:sec:model}).\\
$\Prob(A)$ & The probability of an event $A\subseteq\Omega$: the sum of $p(\omega)$ over the outcomes $\omega\in A$ (Section~\ref{ch13:sec:model}).\\
$\E X$ & The expectation of a random variable $X:\Omega\to\RR$: the weighted average $\sum_{\omega\in\Omega}p(\omega)X(\omega)$ (Section~\ref{ch13:sec:expectation}).\\
$I_A$ & The indicator of the event $A$: it equals $1$ at outcomes in $A$ and $0$ at all other outcomes (Section~\ref{ch13:sec:indicators}).\\
$A_s$, $C_s$ & For a sign sequence $(a_0,\ldots,a_{n-1})$, with every entry $1$ or $-1$: the aperiodic autocorrelation $A_s=\sum_{j=0}^{n-1-s}a_ja_{j+s}$, for $1\le s\le n-1$, with no wrapping around; and the periodic autocorrelation $C_s=\sum_{j=0}^{n-1}a_ja_{(j+s)\bmod n}$, for $0\le s\le n-1$, with positions taken modulo $n$ (Sections~\ref{ch14:sec:aperiodic} and~\ref{ch14:sec:periodic}).\\
\code{a : T}, \code{f x}, \code{:=} & In Lean: \code{a} has type \code{T} (for a proposition \code{P}, \code{h : P} says that \code{h} is a proof of \code{P}); the function \code{f} applied to \code{x}; ``is defined as'' or ``is proved by'' (Section~\ref{ch15:sec:symbols}).\\
\code{sorry} & In Lean, a placeholder for a missing proof. The command \code{\#print axioms} shows whether a theorem depends on one: if it does, the list includes the axiom \code{sorryAx} (Section~\ref{ch16:sec:axioms}).\\
$B(c,r)$ & The ball with center $c$ and radius $r>0$ in a metric space $X$: the points of $X$ whose distance from $c$ is strictly less than $r$ (Section~\ref{ch17:sec:balls}).\\
$y_n$, $\eps_n$, $e_n$ & In Chapter~\ref{ch:capstone}: the points of an iteration whose updates may contain errors, so that $y_{n+1}$ is only close to $T(y_n)$; the update error $\eps_n$, with $d(y_{n+1},T(y_n))\le\eps_n$; and the error $e_n=d(y_n,x_*)$, the distance from $y_n$ to the fixed point (Sections~\ref{ch18:sec:question} and~\ref{ch18:sec:onestep}).\\
\bottomrule
\end{longtable}

\section{Structures, what proofs use, and changes of description}\label{appE:sec:structures}
Almost every object in this book is a set with something extra attached, which Section~\ref{ch06:sec:structure} calls a \emph{structure}. It may be a \emph{relation} (Section~\ref{ch06:sec:differences}), which links or compares pairs of objects, as congruence modulo $m$ and adjacency in a graph do, or an \emph{operation}, such as addition, which combines two elements of a set into an element of the same set. Functions, also called maps (Section~\ref{ch03:sec:functions}), let us carry a question from one description to another.

The extra structure decides which questions can even be asked.
\begin{itemize}
\item Addition and scaling let us form linear combinations of vectors (Chapter~\ref{ch:linear}).
\item Order lets us pick out a smallest or a largest object. Every nonempty set of natural numbers has a smallest member (the well-ordering property of Sections~\ref{ch04:sec:contrapositive} and~\ref{ch05:sec:not-circular}), and every nonempty finite one has a largest; this is how Chapter~\ref{ch:graphs} chooses a shortest walk between two vertices joined by some walk, and a longest directed path in a nonempty finite directed graph.
\item A distance lets us say that one point approximates another, and so speak of limits (Chapter~\ref{ch:limits}). Without one, ``$x_n$ gets close to $x_*$'' has no meaning.
\item A probability on a finite set of outcomes lets us take weighted averages (Chapter~\ref{ch:probability}).
\end{itemize}

Habit~3 (Section~\ref{ch03:sec:images}) asks which information a proof actually uses. For the main theorems of the book the answers are short.
\begin{itemize}
\item The proofs of Hall's theorem (Chapter~\ref{ch:hall}) use only which students accept which tasks and the sizes of finite sets, never what the students and tasks are. So the theorem applies to every finite bipartite graph, whatever its vertices stand for.
\item The proof of the contraction theorem (Theorem~\ref{thm:banach}) uses, besides the facts that the space is nonempty and that $T$ maps it into itself, only the contraction inequality, completeness, and the three requirements of a metric, above all the triangle inequality. It never adds points or multiplies them by numbers, so it works in metric spaces that have no addition at all.
\item The proof of the Cauchy--Schwarz inequality (Theorem~\ref{thm:cs}) uses only three properties of the dot product: linearity in each argument, symmetry, and positivity. So the same proof works for any other way of multiplying two vectors to get a number, provided it has these three properties and the length of $u$ is defined as the square root of the product of $u$ with itself. For example, on $\RR^2$ the product $2u_1v_1+u_2v_2$ is symmetric and linear in each argument, and positive, since $2u_1^2+u_2^2$ is never negative and is zero only when $u=0$. So $|2u_1v_1+u_2v_2|\le\sqrt{2u_1^2+u_2^2}\,\sqrt{2v_1^2+v_2^2}$ for all $u,v\in\RR^2$, with no new work.
\item Chapter~\ref{ch:signs} turns the condition that a sign sequence has zero periodic correlation at every nonzero shift into the condition that different rows of its circulant matrix are orthogonal (Proposition~\ref{ch14:prop:rows}). Adding up all the correlations shows that the order of the matrix, its number of rows, is a perfect square (Section~\ref{ch14:sec:square}), an argument that relies on the circulant pattern. Comparing three rows shows that an order greater than $2$ is divisible by four (Lemma~\ref{thm:hadamardorder}). That argument uses only that the entries are $\pm1$ and different rows are orthogonal, which is exactly what makes a square matrix a Hadamard matrix (Section~\ref{ch14:sec:circulant}), so it holds for every Hadamard matrix, circulant or not.
\end{itemize}
Noticing what a proof does \emph{not} use tells you where its conclusion might reach further. This is a guide for guessing, not a proof: every proposed extension must still be checked against the argument, line by line.

Many steps in the book change the description of a problem without changing the problem, and Habit~6 (Section~\ref{ch06:sec:order}) asks what such a change preserves; a new description helps only when we know exactly what it keeps and what it throws away. Here are four from the book.
\begin{itemize}
\item \emph{An inverse loses nothing.} If $f$ has an inverse $h$, every input is recovered from its output as $x=h(f(x))$. In particular, $f(x)=f(y)$ gives $x=h(f(x))=h(f(y))=y$, so different inputs never share an output, and a question about inputs can be answered by looking at outputs (Section~\ref{ch03:sec:injsurj}).
\item \emph{A quotient forgets on purpose.} Passing from an integer $a$ to its residue class $[a]_m$ forgets everything except the remainder on division by $m$. This is safe only for rules that give the same answer for every representative of a class, which is why well-definedness has to be checked (Section~\ref{ch06:sec:well-defined}).
\item \emph{A matrix turns a linear map into a table of numbers.} Column $j$ of the matrix records where the map sends the standard basis vector $e_j$ (Section~\ref{ch10:sec:linearmaps}), and composing linear maps corresponds to multiplying their matrices (Section~\ref{ch10:sec:product}).
\item \emph{An indicator turns an event into a quantity.} An indicator is $1$ when its event occurs and $0$ when it does not, so counting how many events occur becomes adding indicators, and linearity of expectation applies (Section~\ref{ch13:sec:indicators}).
\end{itemize}
Each of them gives a language in which the step we need is easier to state and to justify.

\section{Choosing the next move}\label{appE:sec:next-move}
Mathematics can look like a long list of tricks to memorize, but it helps far more to recognize the kind of situation you are in. Four situations recur in this book, and each suggests a natural next move.

\subsection*{When a general claim looks doubtful}
Look for a counterexample before you try to prove the claim. Start with small and boundary cases: zero, the empty set, the smallest allowed input, two objects that happen to be equal. A useful counterexample is rarely found by trying cases at random; it is built to break the step that the claim seems to need.
\begin{itemize}
\item Proposition~\ref{thm:compconverse} shows that $g\circ f$ can be injective although $g$ is not. There $f$ maps $\{0,1\}$ into $\{0,1,2\}$ and never reaches $2$, while $g$ sends $2$ to the same output as $1$. The composition applies $g$ only to outputs of $f$, so it never meets this collision.
\item In Section~\ref{ch18:sec:guess}, the guess that small update errors still give convergence to the fixed point fails for a constant bias: the map $x\mapsto qx$ on $\RR$, with $0<q<1$, has the fixed point $0$, but adding the same $\eta>0$ at every update moves the limit to $\eta/(1-q)$, however small $\eta$ is.
\end{itemize}
Each counterexample also showed what a correct statement must take into account: the outputs that $f$ actually reaches, and errors that never die out.

\subsection*{When you need to show that something exists}
When a direct construction is out of reach, ask what a failure would force. The book turns that question into a proof in three ways.
\begin{itemize}
\item \emph{Search and repair.} The search of Chapter~\ref{ch:matching} either finds an augmenting path, which enlarges the matching, or fails and produces a group of students who together accept too few tasks (Proposition~\ref{ch08:prop:shortage}), so that no matching gives every student a task. Either outcome tells us something, which is the point of Habit~7 (Section~\ref{ch08:sec:shortage}).
\item \emph{An extremal choice.} Choose an object that is as large or as small as possible, and show that it already has the property you want. A longest directed path in a nonempty finite acyclic oriented graph ends at a sink, because an arrow leaving its last vertex would give either a longer path or a directed cycle (Theorem~\ref{thm:sink}). Similarly, if a statement about natural numbers fails somewhere, choose the smallest number at which it fails: the statement holds at every smaller number, and this often leads to a contradiction (Section~\ref{ch05:sec:not-circular}).
\item \emph{An average.} A random variable $X$ on a finite probability space satisfies $X(\omega)\ge\E X$ at some outcome $\omega$ of positive probability, since outcomes of probability zero contribute nothing to $\E X$ and a weighted average cannot exceed all the values it averages (Proposition~\ref{ch13:prop:above-average}). Theorem~\ref{thm:cut} uses this: labeling the vertices of a graph with $m$ edges at random gives a cut of expected size $m/2$, so some labeling gives a cut of size at least $m/2$, although the proof never says which. Habit~8 (Section~\ref{ch13:sec:markov}) describes this approach.
\end{itemize}
To choose among them, ask: Can a partial solution be repaired one step at a time? Is there a natural quantity to make as large or as small as possible? Is there an average you can compute exactly?

\subsection*{When the full answer is out of reach}
Ask what you really need to know. Often it is much less than the full answer.
\begin{itemize}
\item To know how close a point $z$ is to the fixed point of a contraction, you need not find the fixed point: the residual estimate of Theorem~\ref{thm:banach} bounds the error by $d(z,Tz)/(1-q)$, and $Tz$ is usually easy to compute. In Section~\ref{ch12:sec:example}, this certified that $x_2=10/27$ lies within $19/729$ of the fixed point of $T(x)=(x^2+1)/3$ on $[0,1]$ before the fixed point had been found.
\item To show that no circulant Hadamard matrix of order $12$ exists, you do not need the complete answer claimed by the research of Section~\ref{ch14:sec:research}. The necessary condition of Proposition~\ref{thm:circulantnecessary} is enough: an order greater than $1$ must have the form $4u^2$ for a positive integer $u$, and $12$ lies strictly between the first two such numbers, $4$ and $16$.
\end{itemize}
The risk in settling for less is answering an easier question than the one you were asked. If you replace the question by an easier one, say so, and compare the new version with the contract of the original, as Habit~2 (Section~\ref{sec:magnitude}) asks. Proposition~\ref{thm:circulantnecessary}, for instance, settles order $12$, but as an answer to ``which orders are possible?'' it is only partial: besides $1$ and $4$, which do occur, it still allows $16,36,64,\ldots$, and says nothing about whether any of them occurs.

\subsection*{When a bound appears}
Ask whether the constant in the bound could be lowered and, if not, which example prevents any improvement. Such an example usually attains the bound, or comes as close to it as we like, so also ask when the bound is attained (Section~\ref{ch09:sec:taste}).
\begin{itemize}
\item The Cauchy--Schwarz inequality $|u\cdot v|\le\norm{u}_2\norm{v}_2$ holds with equality exactly when $v=0$ or $u$ is a multiple of $v$ (Section~\ref{ch10:sec:cs}). Because equality occurs for nonzero vectors, such as $u=v\ne0$, the constant factor $1$ on the right cannot be made smaller.
\item In the constant-bias example of Section~\ref{ch18:sec:sharp}, the error bound of Theorem~\ref{thm:inexact}, part~2, holds with equality at every step: with bias $\eta>0$, the errors are exactly $\eta(1-q^n)/(1-q)$, which approach $\eta/(1-q)$ as $n$ grows. So the factor $1/(1-q)$ is sharp. Any smaller factor would give a bound below the true error once $n$ is large enough, so no cleverer proof can lower it.
\end{itemize}

\section{Turning a reading difficulty into a precise question}\label{appE:sec:question}
Suppose you understand what a theorem says but cannot follow one line of its proof. Asking an assistant, or a classmate, to explain the whole proof again rarely helps: the new explanation may skip the same step, or replace it by a different argument that you must then check from the beginning. It works better to turn the difficulty into a precise question, as Habit~4 (Section~\ref{ch04:sec:promises}) suggests. Write down three things:
\begin{itemize}
\item the exact line that you do not follow;
\item what is known at that point of the proof: the hypotheses, and the facts already established;
\item the last step that you do understand.
\end{itemize}
Then ask which definition or earlier fact leads from what is known to the line. A good answer is usually short, one calculation written out in more detail or one small example, and rarely a new technique. Here are two such questions.
\begin{itemize}
\item In Case~1 of the inductive proof of Hall's theorem (Section~\ref{ch09:sec:inductive}), $S$ is a nonempty proper group of students with $|N(S)|=|S|$, and $T$ is any group of the remaining students. A precise question is: ``I know that the task sets $N(S)$ and $N(T)$ can overlap. Why is $|N(S)\cup N(T)|=|N(S)|+|N(T)\setminus N(S)|$?'' The answer is a splitting: every task in the union lies in exactly one of the sets $N(S)$ and $N(T)\setminus N(S)$, and the sizes of two disjoint finite sets add (Section~\ref{ch03:sec:sets}). Nothing more about graphs is needed. Exercise~\ref{ex:9.3} asks you to write this out in full.
\item From the proof of Theorem~\ref{thm:inexact} (Section~\ref{ch18:sec:onestep}): ``I see that each update error is later multiplied by $q$. Why is the weight of $\eps_j$ in the bound for $e_n$ exactly $q^{n-1-j}$?'' The answer is a count. The error $\eps_j$ is made in the update that produces $y_{j+1}$, and each of the $n-(j+1)=n-1-j$ later updates, which produce $y_{j+2},\ldots,y_n$, contributes one factor $q$.
\end{itemize}

Once you have the answer, write the step in your own words. Then change one small feature and work out what changes.

In the Hall example, the proof continues as follows. The students of $S\cup T$ together accept exactly the tasks in $N(S)\cup N(T)$, so Hall's condition gives $|N(S)\cup N(T)|\ge|S\cup T|$, and $|S\cup T|=|S|+|T|$ because the two groups share no student. The identity from the question, with $|N(S)|=|S|$, turns the left side into $|S|+|N(T)\setminus N(S)|$, and subtracting $|S|$ leaves $|N(T)\setminus N(S)|\ge|T|$. Now suppose $T$ may share students with $S$. The identity still holds, but $|S\cup T|=|S|+|T|$ does not, because a student in both groups is counted once on the left and twice on the right. The conclusion no longer follows, and it can even be false: for $T=S$ it would read $0\ge|S|$, which fails because $S$ is nonempty.

In the second example, if the update errors were numbered from $1$, so that $\eps_j$ bounds the error in the update that produces $y_j$, then $n-j$ updates follow it, and its weight becomes $q^{n-j}$. Or suppose $q=1$, so that the map never increases distances but need not shrink them. If it still has a fixed point $x_*$, the same estimate and count work, but every weight is $1$, so the bound is $e_0+\eps_0+\cdots+\eps_{n-1}$. When every update error equals the same $\eta>0$, this is $e_0+n\eta$, which grows without limit, and the error really can grow this fast: every point is a fixed point of the identity map $T(x)=x$ on $\RR$, so take $x_*=0$, start at $y_0=0$, and set $y_{n+1}=y_n+\eta$; then $e_n=n\eta$, exactly the bound. If you can answer variations like these, the explanation has become your own: you can reason with it, not only repeat it.

The same approach works for a definition that you do not understand. Ask, in this order: What kind of object does it apply to? What exactly must such an object satisfy? What is one example, and what is a \emph{near miss}, an object that fails only one of the requirements? For the definition of an odd integer (Definition~\ref{def:parity}), the objects are integers, and $a$ is odd when $a=2k+1$ for some integer $k$. The number $7=2\cdot3+1$ is an example. The number $8$ is a near miss: $8=2\cdot3.5+1$ has the right shape, but $3.5$, the only $k$ with $2k+1=8$, is not an integer, so the one requirement that fails is ``$k$ is an integer''. Finally, when the definition appears in a proof, decide which of the two directions of Section~\ref{ch01:sec:definition} is used: does the proof \emph{unpack} the definition for an object it is given, or \emph{recognize} that a new object satisfies it? The first proof of the book, Proposition~\ref{thm:odd-sum}, does both: it unpacks the fact that $a$ and $b$ are odd, and recognizes that $a+b$ is even.

\section{What a finished reading record contains}\label{appE:sec:record}
A \emph{reading record} (Section~\ref{ch17:sec:research}) is a written account of a proof you are studying: what you have read, what you have understood, and what remains open. A finished record contains at least:
\begin{itemize}
\item the exact statement of the result, copied from a named source, with the date or version number of what you read, because a preprint or a repository can change;
\item the same statement in your own words;
\item the objects that the statement is about, and the hypotheses on them;
\item a map of the proof's dependencies: a table like the prerequisite map of Section~\ref{ch17:sec:map}, a diagram, or a paragraph;
\item one delicate step of the proof, written out in full detail.
\end{itemize}
Some entries need a note beside them:
\begin{itemize}
\item for a result you use without having read its proof: its hypotheses, and why they hold in your situation;
\item for a computer calculation: the range of cases it covered, and whether it used exact arithmetic or rounded decimals (Section~\ref{appD:sec:fractions});
\item for a result checked by Lean: exactly which statement Lean checked, which axioms its proof depends on (Section~\ref{ch16:sec:axioms}), and whether you ran the check yourself or are reporting what a source says (Section~\ref{ch17:sec:research}).
\end{itemize}

Give every entry the two labels of ``Keeping track of evidence'' in \emph{To the reader}: how well you understand it (\emph{understood}, \emph{can follow}, \emph{borrowed}, or \emph{open}; see also Section~\ref{ch17:sec:map}), and what supports it (proved, cited from another source, checked by computer for finitely many cases, or not established). An entry that rests on one of the four starting facts listed there can say so, since those facts are recorded as assumptions.

The record ends with your own explanation of the proof and one variation you have worked out yourself, such as a hypothesis removed or a conclusion strengthened, with your reasons for what happens. A saved conversation with an assistant cannot replace this record: a conversation mixes guesses, corrections and results without marking which statements survived checking. The record should let another reader rebuild the final argument without reading the conversation.

\section{A suggested capstone rubric}\label{appE:sec:rubric}
The table below suggests one way to assess the independent capstone of Section~\ref{ch18:sec:independent}. Each row names one quality of a finished project, gives its share of the grade, and describes what a strong project shows. The row ``Understanding'' assesses the short oral discussion that ends the investigation. The other rows follow that section's list of what the write-up should contain.
\begin{center}
\begin{tabularx}{\textwidth}{@{}p{1.15in}p{0.5in}X@{}}\toprule
Quality & Weight & What a strong project shows\\\midrule
Claim and scope & 20\% & The exact objects, hypotheses and conclusion of each result. If the question changed during the project, both the original and the revised question are stated, and the difference between them is explained. Each claim is marked with the evidence for it: proved, cited from another source, checked by computer for finitely many cases, or not established.\\
Argument & 35\% & A complete proof or a counterexample, with a map of the background that the argument needs. The delicate steps are written out in full, and each step names the results it depends on.\\
Understanding & 20\% & In the oral discussion, the student explains a delicate step without help, and, when an assumption or an example is changed, correctly predicts what still holds and what fails.\\
Reproducibility & 15\% & Computer checks where they are relevant, with the program, the range of cases it covers and the software used, so that anyone can repeat them. If no computation is relevant, the write-up explains how a reader can check the claims without a computer, for example by working small cases by hand.\\
Context & 10\% & Accurate credit to every source used; a modest statement of what is new compared with the published literature and what is new only to the student; an open account of how an AI assistant was used and of what remains unsettled.\\\bottomrule
\end{tabularx}
\end{center}
Three faults cannot be outweighed by good presentation: a statement that has been changed without saying so, a missing step on which the argument depends, and a claim of novelty without support. A project with any of these faults needs to be corrected before it can pass. The rubric rewards careful work on a small investigation of the student's own; nobody is expected to solve a famous problem.

\section{What you can do at the end of the course}\label{appE:sec:finishing}
You do not need to remember every displayed formula in this book; what matters is what you can do. Faced with a statement, you can say exactly what it requires. Faced with a proof you do not yet follow, you can find where your understanding stops and ask for the one concept or step that would carry you past it. When an explanation arrives, from a book, a classmate or an assistant, you can test it on an example, see its role in the whole proof, and say clearly what has not yet been established.

The next unfamiliar theorem you meet will still contain mathematics you have not seen, but you now have a way in: identify the structures involved, lay out which steps depend on which, learn the missing pieces, and return to the argument. Working with an AI assistant is worthwhile when what it produces helps you understand and gives you something you can check, not when it merely fills a page with text that looks like mathematics.

